\documentclass[12pt,a4paper]{article}

\usepackage[utf8]{inputenc}
\usepackage[T1]{fontenc}
\usepackage{amsmath}
\usepackage{geometry}
\usepackage{booktabs}

\geometry{margin=1in}

\title{Report 1: Gradient Descent}
\author{Dao Quy Tung}
\date{May 9, 2026}

\begin{document}

\maketitle

\section{Introduction}

The objective of this labwork is to implement the gradient descent algorithm from scratch.
Gradient descent is an iterative optimization algorithm used to find a minimum value of a function.
In this report, the tested function is:

\[
f(x) = x^2
\]

Its first-order derivative is:

\[
f'(x) = 2x
\]

The expected minimum value of this function is obtained at:

\[
x = 0, \quad f(x) = 0
\]

\section{Gradient Descent Algorithm}

Gradient descent updates the current value of \(x\) by moving in the opposite direction of the derivative:

\[
x := x - r f'(x)
\]

where:

\begin{itemize}
  \item \(x\) is the current value,
  \item \(r\) is the learning rate,
  \item \(f'(x)\) is the derivative of the function at \(x\).
\end{itemize}

The derivative tells the direction in which the function increases.
Therefore, subtracting the derivative moves \(x\) toward the direction where the function decreases.

\section{Implementation}

The implementation defines two functions: one for \(f(x)\), and one for its derivative \(f'(x)\).
Then, the algorithm repeatedly updates \(x\) until the function value is smaller than a selected threshold.

The main steps of the algorithm are:

\begin{enumerate}
  \item Initialize \(x\), the learning rate \(r\), and the threshold.
  \item Compute the derivative \(f'(x)\).
  \item Update \(x\) using \(x := x - r f'(x)\).
  \item Compute the new value \(f(x)\).
  \item Repeat until \(f(x)\) is small enough.
\end{enumerate}

\section{Experiment Output and Learning Rate Analysis}

The program was executed with:

\[
x_0 = 10, \quad r = 0.9, \quad threshold = 0.01
\]

For this function:

\[
f(x)=x^2, \quad f'(x)=2x
\]

the update rule becomes:

\[
x := x - 0.9(2x)
\]

\[
x := -0.8x
\]

This means the value of \(x\) changes sign after each iteration.
However, its absolute value is multiplied by \(0.8\), so the algorithm still moves closer to zero.

\begin{table}[h]
\centering
\begin{tabular}{cccc}
\toprule
Iteration & \(x\) & \(f'(x)\) & \(f(x)\) \\
\midrule
1 & -8.0000 & -16.0000 & 64.0000 \\
2 & 6.4000 & 12.8000 & 40.9600 \\
3 & -5.1200 & -10.2400 & 26.2144 \\
4 & 4.0960 & 8.1920 & 16.7772 \\
5 & -3.2768 & -6.5536 & 10.7374 \\
6 & 2.6214 & 5.2429 & 6.8719 \\
7 & -2.0972 & -4.1943 & 4.3980 \\
8 & 1.6777 & 3.3554 & 2.8147 \\
9 & -1.3422 & -2.6844 & 1.8014 \\
10 & 1.0737 & 2.1475 & 1.1529 \\
11 & -0.8590 & -1.7180 & 0.7379 \\
12 & 0.6872 & 1.3744 & 0.4722 \\
13 & -0.5498 & -1.0995 & 0.3022 \\
14 & 0.4398 & 0.8796 & 0.1934 \\
15 & -0.3518 & -0.7037 & 0.1238 \\
16 & 0.2815 & 0.5629 & 0.0792 \\
17 & -0.2252 & -0.4504 & 0.0507 \\
18 & 0.1801 & 0.3603 & 0.0325 \\
19 & -0.1441 & -0.2882 & 0.0208 \\
20 & 0.1153 & 0.2306 & 0.0133 \\
21 & -0.0922 & -0.1845 & 0.0085 \\
\bottomrule
\end{tabular}
\caption{Output of the implemented gradient descent program.}
\end{table}

The output shows that \(x\) oscillates around zero because the learning rate is large.
Nevertheless, \(f(x)\) decreases from \(64.0000\) to \(0.0085\), which is below the threshold \(0.01\).
Therefore, the algorithm stops after 21 iterations.

The learning rate controls how large each update step is.
For the function:

\[
f(x) = x^2
\]

the update rule is:

\[
x := x - r(2x)
\]

or:

\[
x := (1 - 2r)x
\]

This shows that the value of \(r\) strongly affects convergence.

\subsection{Small Learning Rate}

If \(r\) is too small, the algorithm moves slowly.
The loss decreases, but many iterations are needed before reaching the minimum.

For example, if \(r = 0.01\):

\[
x := 0.98x
\]

Only a small part of \(x\) is reduced at each step.

\subsection{Suitable Learning Rate}

If \(r\) is suitable, the algorithm converges quickly and stably.
For example, if \(r = 0.1\):

\[
x := 0.8x
\]

This makes \(x\) decrease smoothly toward zero.

\subsection{Large Learning Rate}

If \(r\) is too large, the algorithm may overshoot the minimum.
For example, if \(r = 0.9\):

\[
x := (1 - 2 \times 0.9)x = -0.8x
\]

The sign of \(x\) changes at each iteration.
However, because the absolute value is multiplied by \(0.8\), the algorithm still converges, but it oscillates around zero.

If \(r\) is even larger, the algorithm can diverge.
For example, if \(r = 1.1\):

\[
x := (1 - 2 \times 1.1)x = -1.2x
\]

The absolute value of \(x\) becomes larger after each iteration, so the algorithm moves away from the minimum.

\section{Conclusion}

In this labwork, gradient descent was implemented from scratch and tested with the function \(f(x)=x^2\).
The derivative \(f'(x)=2x\) provides the direction of increase, so the algorithm subtracts the derivative multiplied by the learning rate to reduce the function value.
The experiment shows that the learning rate is very important: a small learning rate makes convergence slow, a suitable learning rate gives stable convergence, and a very large learning rate can cause oscillation or divergence.

\end{document}
